% !TITLE 八声甘州·对潇潇暮雨洒江天
% !AUTHOR 柳永
% !DYNASTY 两宋
% !GENRE 词
% !ORDER 5

\begin{poem}
对潇潇\textsuperscript{1}暮雨洒江天，一番洗清秋。
渐霜风凄紧\textsuperscript{2}，关河\textsuperscript{3}冷落，残照当楼。
是处红衰翠减\textsuperscript{4}，苒苒物华休\textsuperscript{5}。
唯有长江水，无语东流。

不忍登高临远，望故乡渺邈\textsuperscript{6}，归思难收。
叹年来踪迹，何事苦淹留\textsuperscript{7}？
想佳人、妆楼颙望\textsuperscript{8}，误几回、天际识归舟\textsuperscript{9}。
争知我\textsuperscript{10}、倚栏杆处，正恁凝愁\textsuperscript{11}！
\end{poem}

\begin{annotations}
\notetext{1}{潇潇：风雨声，也形容雨势急骤。傍晚的雨声洒在江天之间。}
\notetext{2}{霜风凄紧：秋风刮得凄厉紧迫。霜风，寒冷的风。}
\notetext{3}{关河：关山与河流，泛指山河大地。}
\notetext{4}{是处红衰翠减：到处红花凋谢、绿叶枯萎。是处，到处；红，代花；翠，代叶。}
\notetext{5}{苒苒物华休：美好的景物渐渐消逝。苒苒（rǎn），渐渐；物华，自然美景；休，结束。}
\notetext{6}{渺邈：遥远模糊。渺邈（miǎo miǎo），极远而看不清。}
\notetext{7}{淹留：长久滞留他乡。为什么苦苦地久留在这异乡呢？}
\notetext{8}{颙望：抬头远望。颙（yóng），仰头。想象佳人在妆楼上抬头望向远方，盼着归船。}
\notetext{9}{天际识归舟：在天边辨认归来的船只。多少次以为看到了归船，走近才发现认错了——写尽了等待中的希望和失望。}
\notetext{10}{争知我：怎么知道我现在……争，怎么（宋代口语）。}
\notetext{11}{恁凝愁：如此忧愁凝结不解。恁（nèn），那样、如此；凝愁，愁绪郁结凝固。}
\end{annotations}

\begin{appreciation}
《八声甘州》是柳永的代表作之一，写的是游子漂泊在外、想念家乡。苏轼夸它“不减唐人高处”——能让苏轼低头认账的词不多。

“对潇潇暮雨洒江天，一番洗清秋。”傍晚的雨洒满江面、洒满天空，把秋天洗得干干净净。一个“洗”字，秋天亮了。“渐霜风凄紧，关河冷落，残照当楼”——一个“渐”字把时间慢慢拉开：雨停，风起，山河冷落，夕阳斜照在楼上。然后是“是处红衰翠减，苒苒物华休”：到处红花谢了、绿叶枯了，好风景一样一样地结束。这和李清照的“绿肥红瘦”是同一件事——她看的是花园，柳永看的是万里江山。

“唯有长江水，无语东流。”只有长江水不在乎季节，年年无语东流。人看着流水，就想起“年来踪迹”——一年年在外头飘，到底图什么。“何事苦淹留”，他自己也回答不上来。

下阕写思乡。“不忍登高临远，望故乡渺邈，归思难收。”不敢登高——登高了也看不见故乡，可思乡的心收不住。最妙的是接下来：“想佳人、妆楼颙望，误几回、天际识归舟。”他不写自己想家，他写家里的她：她站在妆楼上抬头望，一次一次把天边的船认成他的船，又一次一次认错。读《行行重行行》时，我们看见的是等人的女子；这里柳永替走的人，把等人的心也写出来了。同一场离别，两头的苦，他都尝到了。

结尾一句，“争知我、倚栏杆处，正恁凝愁”——她也不知道，我正靠着栏杆，愁得像化不开的夜色。两个人互相想，却互相不知道，这比一个人想更寂寞。

柳永是北宋第一个专心写词的文人。他考不上科举，一辈子在江湖上飘，所以写漂泊写得比谁都真——他的词当时传唱得满城都是。

以后你离开家去读书、去工作，多半也会有几个想家的晚上。想家的时候给家里打个电话，比“误几回、天际识归舟”强。别一个人把愁凝在那里。
\end{appreciation}
